\chapter{可控知识编辑内容安全防御系统设计与开发}\label{chap:experiment}

\section{需求分析}
本节为示例占位内容。表标题和表序号居中置于表上方，如表~\ref{tab:params} 所示。

\begin{table}[htbp]
  \centering
  \caption{示例表格——主要参数设置}
  \label{tab:params}
  \begin{tabular}{lcc}
    \toprule
    参数 & 符号 & 取值 \\
    \midrule
    参数一 & $\alpha$ & $0.1$ \\
    参数二 & $\beta$  & $0.9$ \\
    参数三 & $N$      & $1000$ \\
    \bottomrule
  \end{tabular}
\end{table}

\subsection{用户分析}
% TODO: 安全编辑从实验脚本到可用平台的工程缺口；目标用户与使用场景。

\subsection{功能需求}
% TODO: 编辑管理、分级配置、持久性监控、安全评估四类功能需求。

\subsection{非功能需求}
% TODO: 性能、可扩展性（分级映射可替换）、评估口径统一等非功能需求。

\section{系统概要介绍}

总体架构设计
% TODO: 编辑引擎层 / 分级配置层 / 持久性监控层 / 安全评估层四层架构。

技术选型
% TODO: EasyEdit 封装、trl 微调、评估流水线、前端框架与开发环境。

系统工作流程
% TODO: 编辑配置 → 执行编辑 → 微调模拟 → 持久性监控 → 安全评估的端到端流程。

\section{功能模块设计}
本节为示例占位内容。不同方法的对比结果如表~\ref{tab:results} 所示，其中表头
使用了跨行单元格。

\begin{table}[htbp]
  \centering
  \caption{示例表格——结果对比}
  \label{tab:results}
  \begin{tabular}{lccc}
    \toprule
    \multirow{2}{*}{方法} & 指标一 & 指标二 & 指标三 \\
                          & (\%)   & (\%)   & (s)    \\
    \midrule
    方法 A     & $86.0$ & $5$ & $42.3$ \\
    方法 B     & $91.5$ & $3$ & $38.7$ \\
    本文方法   & $\mathbf{98.2}$ & $\mathbf{0}$ & $\mathbf{31.5}$ \\
    \bottomrule
  \end{tabular}
\end{table}

由定理~\ref{thm:demo} 与式~\eqref{eq:demo1} 可知，此处可进一步展开对结果的分析说明。

\subsection{编辑管理模块}
% TODO: 编辑方法与强度选择、编辑任务队列与权重版本管理（创新点一方法落地）。

\subsection{分级配置模块}
% TODO: LEVEL_MAP 可视化配置、风险分级器接入（创新点二方法落地）。

\subsection{持久监控模块}
% TODO: 微调模拟与衰减曲线 dashboard，编辑保持率实时监控。

\subsection{安全评估模块}
% TODO: 攻击成功率与过度拒绝率自动化评估流水线，统一判定口径。

\section{系统开发}

核心模块实现
% TODO: 四大模块的实现要点与关键代码说明。

关键技术难点
% TODO: 模型权重管理、评估并发、判定口径统一等关键问题的解决方案。

\section{系统测试}

功能测试
% TODO: 各模块功能测试用例与测试结果。

集成效果验证
% TODO: 集成第 2/3 章方法的端到端效果验证（安全编辑 → 微调模拟 → 评估全流程）。

性能测试
% TODO: 编辑耗时、评估吞吐、监控延迟等性能指标测试。

\section{本章小结}
本节为示例占位内容。请替换为本章的小结。
